\documentclass[11pt]{article}
\usepackage[margin=2.7cm]{geometry}
\usepackage{amsmath,amssymb,amsthm}
\usepackage{booktabs}
\usepackage{xcolor}
\definecolor{linkblue}{rgb}{0.1,0.1,0.45}
\usepackage[colorlinks=true,linkcolor=linkblue,citecolor=linkblue,urlcolor=linkblue]{hyperref}
\usepackage{microtype}

\newtheorem{theorem}{Theorem}[section]
\newtheorem{lemma}[theorem]{Lemma}
\newtheorem{corollary}[theorem]{Corollary}
\newtheorem{conjecture}[theorem]{Conjecture}
\theoremstyle{definition}
\newtheorem{definition}[theorem]{Definition}
\newtheorem{hypothesis}[theorem]{Hypothesis}
\newtheorem{example}[theorem]{Example}
\theoremstyle{remark}
\newtheorem{remark}[theorem]{Remark}

\newcommand{\A}{\mathcal{A}}
\newcommand{\Hs}{\mathcal{H}}
\newcommand{\Z}{\mathbb{Z}}
\newcommand{\R}{\mathbb{R}}
\newcommand{\C}{\mathbb{C}}
\newcommand{\T}{\mathbb{T}}
\newcommand{\tr}{\operatorname{tr}}
\newcommand{\supp}{\operatorname{supp}}
\newcommand{\spec}{\operatorname{spec}}
\newcommand{\slim}{\operatorname*{s-lim}}

\title{Quantization of the domain-wall memory effect\\ on a spin chain, without scattering assumptions\\[1ex]
\large A self-contained report}
\author{Project \texttt{tns-as-st-me} --- campaign report\thanks{Prepared 2026-08-28 by the
project's automated adversarial-verification campaign.  Every theorem below was
produced by a proposer--critic loop in which independent provers and hostile
reviewers (from different model families) attacked each statement until no
fatal or major objection remained; the final adjudication records are cited in
Appendix~\ref{app:method}.  The internal claim identifiers
(\texttt{M-INDEX-fin}, \texttt{M-INDEX-spec}, \texttt{M-IDX-density},
\texttt{AC-EX}) refer to the project's claim ledger.}}
\date{28 August 2026}

\begin{document}
\maketitle

\begin{abstract}
\noindent
A wall between two oppositely magnetized halves of a spin chain shifts its
position when a spin wave scatters off it, and the shift survives after the
spin wave has left: the chain remembers the collision.  This report proves,
in a precisely delimited sense, that this remembered shift is \emph{quantized}
--- each recorded outcome is an integer multiple of a unit fixed by the
magnetization density --- and that the proof needs \emph{no} assumptions about
how the scattering proceeds: no decomposition into reflected and transmitted
waves, no absence of bound states, no completeness of any scattering theory.
Those scattering assumptions, previously believed necessary, are needed only
to compute \emph{which} multiple occurs with which probability, not to prove
that only integer multiples occur.  Three further results are established
along the way: the average shift itself is \emph{not} quantized (it is a
probability-weighted average of quantized outcomes, and the distinction is
essential); a by-product theorem shows that the magnetization density of a
symmetric matrix-product ground state is itself quantized in half-integers;
and there is provably \emph{no} single operator whose spectrum expresses the
quantization --- the measurement protocol in the statement cannot be removed.
Finally, for chains whose excitations are described exactly by the
matrix-product excitation ansatz, the existence half of the scattering theory
is proved from the ansatz data.  Every term used above is defined in the body
of the report.  Each theorem's exact hypotheses, including two assumptions
that remain unproved from first principles, are stated in full; a section
lists what has \emph{not} been proved.  All results were verified by
adversarial review and by machine-checked certificates, including a
pre-registered numerical falsification attempt that failed to falsify.
\end{abstract}

\tableofcontents

\section{Summary of results}\label{sec:summary}

The whole content of this report in one sentence: \emph{on an infinite spin
chain carrying a wall between two domains of opposite magnetization density
$\pm\rho$, the outcome of measuring the wall's displacement --- by the
explicit two-measurement protocol of Definition~\ref{def:tpm} --- is always an
integer multiple of $1/(2\rho)$ lattice sites, provided only that the chain's
local dynamics conserves magnetization and that the measured region settles
down at late times (Hypothesis~\ref{hyp:lr}); no property of the scattering
process is assumed.}

The five main statements, with their logical status:

\begin{center}
\begin{tabular}{@{}llp{7.8cm}@{}}
\toprule
 & status & content \\
\midrule
Theorem~\ref{thm:fin} & proved, unconditional & at any fixed finite window,
the measured charge lies on a rigid integer lattice, at all times \\
Theorem~\ref{thm:spec} & proved from Hyp.~\ref{hyp:int}+\ref{hyp:lr} &
the asymptotic outcome distribution of the two-measurement protocol is
supported on the integers; the mean shift is
$\delta x = -(2\rho)^{-1}\sum_\nu \nu\, p_\nu$ \\
Theorem~\ref{thm:density} & proved, unconditional & the density itself obeys
$2\rho\in\Z$ for symmetric injective matrix-product vacua \\
Theorem~\ref{thm:noq} & proved (a refutation) & no global charge operator
exists whose spectrum could express Theorem~\ref{thm:spec}; the protocol is
unremovable \\
Theorem~\ref{thm:acex} & proved from Hyp.~\ref{hyp:band} & scattering states
(wave operators) \emph{exist} whenever the matrix-product excitation ansatz is
exact; their completeness is not claimed and not needed \\
\bottomrule
\end{tabular}
\end{center}

Section~\ref{sec:notproved} lists, without euphemism, what is assumed and what
remains open.  Section~\ref{sec:numerics} reports the numerical checks,
including a falsification attempt that was specified before the theorems were
proved.

\section{The setting}\label{sec:setting}

Everything in this report happens on an infinite one-dimensional lattice of
quantum spins.  This section defines the lattice, its observables, its states,
and its dynamics.  A reader familiar with quantum spin systems can skim to
Section~\ref{sec:mps}; nothing nonstandard appears before then.

\begin{definition}[spin chain, local observables]\label{def:chain}
Fix an integer $d\ge 2$ (the \emph{on-site dimension}).  To each site
$x\in\Z$ attach the Hilbert space $\C^d$.  For a finite set
$\Lambda\subset\Z$, an \emph{observable supported in $\Lambda$} is a linear
operator on $\bigotimes_{x\in\Lambda}\C^d$, extended to act as the identity on
all other sites.  The collection of all such operators, for all finite
$\Lambda$, together with the limits of sequences that converge in operator
norm, is denoted $\A$ and called the \emph{observable algebra} of the chain.
\end{definition}

\begin{definition}[state]\label{def:state}
A \emph{state} $\omega$ is a linear map $\omega:\A\to\C$ assigning to each
observable its expectation value, with $\omega(A^\dagger A)\ge 0$ for all $A$
and $\omega(\mathbf{1})=1$.  This is the standard notion: a state is a
consistent assignment of expectation values to all local measurements.
\end{definition}

\begin{definition}[Hilbert-space realization]\label{def:gns}
Every state $\omega$ can be \emph{realized} on a Hilbert space: there exist a
Hilbert space $\Hs$, a unit vector $\Psi\in\Hs$, and an assignment $\pi$ of an
operator $\pi(A)$ on $\Hs$ to each $A\in\A$, compatible with sums, products
and adjoints, such that $\omega(A)=\langle\Psi,\pi(A)\Psi\rangle$ for all
$A\in\A$.  (This is the Gelfand--Naimark--Segal construction; we use only its
existence.)  When a state is fixed we write $A$ for $\pi(A)$.
\end{definition}

\begin{definition}[on-site charge, total charge of a region]\label{def:charge}
Fix a Hermitian $d\times d$ matrix $S^z$ (the \emph{on-site charge}; for a
spin-$S$ chain, $d=2S+1$ and $S^z$ is the usual third spin component with
eigenvalues $-S,-S+1,\dots,S$).  Write $S^z_x$ for the copy of $S^z$ at site
$x$.  For a finite interval $W=[a,b]\subset\Z$ the \emph{charge of $W$} is
$\sum_{x=a}^{b}S^z_x$.  The one-parameter family of unitaries
$u(\theta)=e^{i\theta S^z}$, applied simultaneously at every site of a region,
implements a \emph{rotation} of that region by angle $\theta$.
\end{definition}

\begin{hypothesis}[integer spacing of the on-site charge, (INT)]\label{hyp:int}
There is a complex number $c$ with $|c|=1$ such that
$e^{2\pi i S^z}=c\,\mathbf{1}$ on one site.  Equivalently, writing
$c=e^{2\pi i\kappa}$ with $\kappa\in[0,1)$: all eigenvalues of $S^z$ lie in
the shifted integer lattice $\kappa+\Z$.
\end{hypothesis}

\begin{remark}
Hypothesis~\ref{hyp:int} holds for every spin-$S$ chain, with
$c=e^{2\pi iS}$: the eigenvalues of the third spin component are spaced by
exactly $1$.  It is the only structural input about the on-site charge that
any result below uses.  No result assumes anything about the \emph{value} of
any density; densities are outputs (Theorem~\ref{thm:density}), not inputs.
\end{remark}

\begin{definition}[dynamics]\label{def:dynamics}
A \emph{finite-range, charge-conserving Hamiltonian} is a formal sum
$H=\sum_{x\in\Z}h_x$ where each $h_x$ is a fixed Hermitian observable
supported in $[x,x+R]$ for some fixed range $R<\infty$, independent of $x$ up
to translation, and each $h_x$ commutes with the total charge of any interval
containing its support plus a collar of width $R$.  Such a sum generates, in a
standard way, a family of transformations $\alpha_t$ ($t\in\R$) of the
observable algebra: $\alpha_t(A)$ is the observable $A$ evolved for time $t$
(Heisenberg picture).  Charge conservation means concretely: the change of the
charge of $W$ in time equals the net current through the two ends of $W$,
where the \emph{current} at a bond is itself a local observable determined by
$H$; no charge is created or destroyed in the interior.
\end{definition}

\section{The two vacua, matrix product states, and the density theorem}
\label{sec:mps}

The wall of the title separates two translation-invariant reference states.
This section defines the class of reference states used --- matrix product
states --- and proves the first main theorem, which constrains their
densities.  The definitions here are standard in the tensor-network
literature \cite{cirac2021}; they are restated so that this report is
self-contained.

\begin{definition}[matrix product state, MPS]\label{def:mps}
Fix an integer $\chi\ge 1$ (the \emph{bond dimension}) and a collection
$A=(A^{1},\dots,A^{d})$ of $\chi\times\chi$ complex matrices, one per basis
state of the site Hilbert space $\C^d$.  Define the linear map
$\mathbb{E}$ on $\chi\times\chi$ matrices, and for each single-site operator
$X$ the map $\mathbb{E}_X$, by
\[
\mathbb{E}(Y):=\sum_{m=1}^{d}A^{m}\,Y\,(A^{m})^{\dagger},\qquad
\mathbb{E}_X(Y):=\sum_{m,n=1}^{d}\langle n|X|m\rangle\,A^{m}\,Y\,(A^{n})^{\dagger}.
\]
Assume the normalization: the largest eigenvalue of $\mathbb{E}$ (in absolute
value) is $1$, it is the unique eigenvalue on the unit circle, and it has
positive-definite fixed points --- a matrix $\Lambda$ with
$\mathbb{E}(\Lambda)=\Lambda$ and a matrix $\Lambda'$ with
$\mathbb{E}^{\dagger}(\Lambda')=\Lambda'$, normalized by
$\tr(\Lambda'\Lambda)=1$.  The \emph{matrix product state} $\omega_A$ is the
translation-invariant state whose expectation of a product
$X_{1}\otimes\cdots\otimes X_{n}$ of single-site operators on $n$ consecutive
sites is
\[
\omega_A\bigl(X_{1}\otimes\cdots\otimes X_{n}\bigr)
=\tr\bigl(\Lambda'\;\mathbb{E}_{X_{1}}\circ\cdots\circ\mathbb{E}_{X_{n}}(\Lambda)\bigr).
\]
The matrices $A^m$ act on an auxiliary $\chi$-dimensional space (the
\emph{bond space}) that carries the correlations of the state; $\chi=1$ gives
exactly the uncorrelated product states.
\end{definition}

\begin{definition}[injectivity]\label{def:injective}
The tensor $A$ is \emph{injective} if there is a length $k$ such that the
products $A^{m_1}A^{m_2}\cdots A^{m_k}$, over all choices of indices, span the
full space of $\chi\times\chi$ matrices.  Injectivity is the genericity
condition of the MPS literature: it guarantees that the state determines the
tensor up to the natural equivalences, and that the normalization of
Definition~\ref{def:mps} can be arranged.  All reference states below are
injective MPS.
\end{definition}

\begin{definition}[symmetric MPS pair with opposite densities]\label{def:vacua}
Fix two injective MPS tensors $A_\alpha$, $A_\beta$ of a common bond dimension
$\chi$, with states $\omega_\alpha$, $\omega_\beta$ (the two \emph{vacua}).
We assume:
\begin{enumerate}
\item \emph{(invariance)} both states are invariant under all rotations:
$\omega_\gamma\bigl(u(\theta)^{\otimes n}\,B\,(u(\theta)^{\dagger})^{\otimes n}\bigr)
=\omega_\gamma(B)$ for $\gamma\in\{\alpha,\beta\}$, every $\theta$, and every
observable $B$ on $n$ sites;
\item \emph{(local implementation)} for each vacuum there are invertible
$\chi\times\chi$ matrices $V_\theta$ and phases $f(\theta)$, with
$V_0=\mathbf{1}$, $f(0)=0$, both differentiable in $\theta$, such that
rotating the physical index of the tensor equals a similarity transformation
on the bond space:
\[
\sum_{n}\,[u(\theta)]_{mn}\,A^{n}
=e^{i f(\theta)}\;V_\theta\,A^{m}\,V_\theta^{-1}\qquad(m=1,\dots,d).
\]
For an injective tensor, clause (2) with \emph{some} $V_\theta$ and
$f(\theta)$ is a theorem given clause (1) (the ``fundamental theorem'' of MPS
symmetries \cite{perez2008,cirac2021}); the differentiability of the chosen
family is the one extra smoothness assumption, labelled (S);
\item \emph{(opposite densities)} the two charge densities are antisymmetric:
\[
\rho\;:=\;\omega_\alpha(S^z_0)\;=\;-\,\omega_\beta(S^z_0).
\]
Here $\rho$ is a real parameter of the pair; nothing is assumed about its
value, and $\rho=0$ is allowed;
\item \emph{(stationarity)} both vacua are invariant in time under the
dynamics of Definition~\ref{def:dynamics}, and $H$ commutes with rotations.
\end{enumerate}
\end{definition}

\begin{example}
The simplest example is the spin-$\tfrac12$ chain with
$\alpha=$ all spins up, $\beta=$ all spins down: product states ($\chi=1$)
with $\rho=\tfrac12$.  A genuinely correlated example with $\rho=0$ is the
spin-$1$ chain state of Affleck, Kennedy, Lieb and Tasaki (the \emph{AKLT
state}), an explicit injective MPS with $\chi=2$.  The point of the definitions above is that
\emph{nothing} below is restricted to product states.
\end{example}

The first theorem was found as a by-product of the campaign and constrains
which densities can occur at all.  It has the flavor of the
Lieb--Schultz--Mattis-type results, which tie densities of symmetric ground
states to lattice constraints, but its proof is a short exact computation.

\begin{theorem}[density quantization; claim \texttt{M-IDX-density}]
\label{thm:density}
Under Definition~\ref{def:vacua} (clauses 1--3, with the smoothness (S) at
both vacua) and Hypothesis~\ref{hyp:int} with constant $c$:
\[
e^{2\pi i\rho}\;=\;c\;=\;e^{-2\pi i\rho},
\qquad\text{and therefore}\qquad 2\rho\in\Z .
\]
\end{theorem}

\begin{proof}[Proof sketch]
Differentiating the local-implementation identity of
Definition~\ref{def:vacua}(2) at $\theta=0$ and taking the expectation value
in $\omega_\alpha$ gives $f'(0)=\omega_\alpha(S^z_0)=\rho$: the phase slope
\emph{is} the density.  The identity composes under successive rotations, and
injectivity makes the decomposition into ($f$, $V_\theta$) rigid enough that
$f$ is additive; with differentiability this forces
$f(\theta)=\rho\,\theta$ up to integer multiples of $2\pi$.  At $\theta=2\pi$
the rotation $u(2\pi)=c\,\mathbf{1}$ is a pure number by
Hypothesis~\ref{hyp:int}, so $e^{2\pi i\rho}=c$.  Running the same argument in
the second vacuum, whose density is $-\rho$, gives $e^{-2\pi i\rho}=c$.
Dividing the two conclusions yields $e^{4\pi i\rho}=1$.  The antisymmetry of
clause (3) is load-bearing: one vacuum alone gives only $\rho\in\kappa+\Z$,
and a general pair only $\rho_\alpha-\rho_\beta\in\Z$.  No step uses any
relation between $\rho$ and the on-site dimension $d$.  The full proof, in
machine-auditable numbered steps, is in the project file
\texttt{theory/memory-index.md}; its two load-bearing computations carry
machine-checked certificates (Section~\ref{sec:numerics}).
\end{proof}

\section{The wall, the window charge, and the wall position}
\label{sec:wall}

\begin{definition}[kink states]\label{def:kink}
A \emph{kink state} (or wall state) is a state $\omega$ of the chain that
agrees with $\omega_\alpha$ far to the left and with $\omega_\beta$ far to the
right, with summable charge deviations:
\[
\sum_{x<0}\bigl|\omega(S^z_x)-\rho\bigr|
\;+\;\sum_{x>0}\bigl|\omega(S^z_x)+\rho\bigr| \;<\;\infty .
\]
When first moments are needed (they are, in Hypothesis~\ref{hyp:lr}(iii)) we
additionally require the sums above to converge with an extra weight $|x|$,
i.e.\ $\sum_x |x|\,\bigl|\omega(S^z_x)\mp\rho\bigr|<\infty$ with the sign
matching the corresponding side.  From here on, fix a kink
state, realize it on a Hilbert space as in Definition~\ref{def:gns}, and call
the realizing vector $\Psi$.
\end{definition}

For $\rho>0$ the charge profile of a kink state falls from $+\rho$ on the far
left to $-\rho$ on the far right; the crossover region is the \emph{wall}.
Moving the entire wall one site to the right converts one site's expected
charge from $-\rho$ to $+\rho$: the region of density $+\rho$ grows by one
site, so the total charge \emph{increases} by $2\rho$.  This calibration --- one site of wall motion equals $2\rho$ units of
charge --- is the entire mechanism of the memory effect, and motivates:

\begin{definition}[window charge and wall position]\label{def:window}
For a finite interval $W=[a,b]$ and a marked interior site $c\in W$ (the
\emph{cut}), the \emph{window charge} is the observable
\[
\widehat{Q}_{W,c}\;:=\;\sum_{x=a}^{c}\bigl(S^z_x-\rho\bigr)
\;+\;\sum_{x=c+1}^{b}\bigl(S^z_x+\rho\bigr),
\]
i.e.\ the charge of $W$ measured against the profile of a reference wall
sitting exactly at the cut.  The associated \emph{wall-position observable}
is
\[
X_{W,c}\;:=\;c\;+\;\frac{\widehat{Q}_{W,c}}{2\rho}\qquad(\rho>0),
\]
so that a state whose wall sits $n$ sites to the right of $c$ (and which is
exactly vacuum elsewhere in $W$) has $\langle X_{W,c}\rangle=c+n$.
\end{definition}

\begin{remark}
$\widehat{Q}_{W,c}$ is a bounded observable supported in the finite window
$W$; no infinite sums are involved.  All limits ($t\to\pm\infty$ first, then
$W\uparrow\Z$) are taken later, in a stated order, and the order matters.
\end{remark}

\section{What ``memory'' means, and why the naive definition fails}
\label{sec:naive}

The physical effect: prepare a kink plus an incoming localized spin-wave
packet (a \emph{magnon}: a traveling disturbance of the vacuum carrying one
unit of charge); let them collide; after the collision the wall sits, on
average, somewhere else.  The displacement persists after the magnon has
left --- the chain remembers.  The question is what exactly is quantized.

One is tempted to define a displacement \emph{operator}
``$X(+\infty)-X(-\infty)$'' and ask for its spectrum.  Two theorems below say
this temptation must be resisted, for two independent reasons.

\emph{First reason: differences of lattice-valued observables are not
lattice-valued.}  The window charge at late times,
$\alpha_t(\widehat{Q}_{W,c})$, and at early times,
$\alpha_{-t}(\widehat{Q}_{W,c})$, do not commute with one another.  For
noncommuting operators, even when each has spectrum in $\Z$, the difference
need not: take, on $\C^2$ with orthonormal basis $\{|0\rangle,|1\rangle\}$ and
$|+\rangle:=(|0\rangle+|1\rangle)/\sqrt2$, the projections
$Q_-=|0\rangle\langle 0|$ and $Q_+=|+\rangle\langle +|$.  Both have spectrum
$\{0,1\}\subset\Z$, yet
\[
\spec\,(Q_+-Q_-)=\Bigl\{-\tfrac{1}{\sqrt2},\,+\tfrac{1}{\sqrt2}\Bigr\},
\]
which contains no integers at all.  Any correct quantization statement must
therefore say \emph{how} the two charges are measured, not merely subtract
them.

\emph{Second reason: the global charge operator does not exist.}  One might
instead try to build a single conserved operator (``total charge relative to
the wall'') and use its spectrum.  This fails:

\begin{theorem}[no global charge operator; claim \texttt{M-INDEX-LA-strong},
refuted form]\label{thm:noq}
There is no construction of a self-adjoint operator $\widehat{Q}$ on the
Hilbert space of Definition~\ref{def:kink} that (i) is the limit of the
window charges $\widehat{Q}_{W,c}$ as $W\uparrow\Z$ in any of the standard
operator senses, and (ii) exists for every kink state satisfying the
summability condition of Definition~\ref{def:kink}.  Two independent
obstructions were established:
\begin{enumerate}
\item \emph{(fluctuations)} the summability condition controls the
\emph{means} $\omega(S^z_x)\mp\rho$ but not the fluctuations; there are
explicit product states (bond dimension $1$, zero Hamiltonian) satisfying
Definition~\ref{def:kink} for which the variance of $\widehat{Q}_{W,c}$
diverges logarithmically as $W$ grows, so the sequence converges to no
operator;
\item \emph{(correlated vacua)} whenever the bond matrices $V_\theta$ of
Definition~\ref{def:vacua}(2) are not multiples of the identity --- the
generic situation for $\chi\ge2$ --- the half-infinite partial sums of
$(S^z_x-\rho)$ fail to converge on the kink vector, by the same mechanism
that makes a rotation of half the chain impossible to implement by a
norm-convergent sequence of local unitaries.
\end{enumerate}
Consequently any quantization theorem must be a statement about
\emph{measurement outcomes}, not about the spectrum of a global observable.
A weaker existence statement --- a charge operator built from the symmetry
group acting on the specific Hilbert space realization of the kink state ---
remains plausible and is recorded as a conjecture in
Section~\ref{sec:notproved}.
\end{theorem}

\section{The two-measurement protocol and the relaxation hypothesis}
\label{sec:protocol}

The correct object is the outcome statistics of an explicit, physically
performable protocol.

\begin{definition}[projective measurement of the window charge]
\label{def:projmeas}
By Theorem~\ref{thm:fin} below, $\widehat{Q}_{W,c}$ has a discrete set of
eigenvalues.  \emph{Measuring $\widehat{Q}_{W,c}$ at time $t$} means the
textbook projective measurement of the Heisenberg observable
$\alpha_t(\widehat{Q}_{W,c})$: the outcome is an eigenvalue $q$, obtained
with probability $\|E_{W,t}(\{q\})\Psi\|^2$, where $E_{W,t}$ denotes the
projection onto the eigenspace of $\alpha_t(\widehat{Q}_{W,c})$ with
eigenvalue $q$; after the measurement the state is the normalized projection.
\end{definition}

\begin{definition}[the two-measurement protocol and its outcome law]
\label{def:tpm}
Fix the window $W$, the cut $c$, an early time $-t_-$ and a late time $t_+$
(both large).  The protocol: (i) measure $\widehat{Q}_{W,c}$ at time $-t_-$,
record the outcome $q_-$; (ii) let the system evolve; (iii) measure
$\widehat{Q}_{W,c}$ at time $t_+$, record $q_+$; (iv) record the
\emph{escaped charge}
\[
\nu\;:=\;q_- - q_+ ,
\]
the amount of charge that left the window between the two measurements.  The
\emph{outcome law} $p_{W;t_-,t_+}(\nu)$ is the probability of recording the
value $\nu$, computed by the standard rules of sequential projective
measurement:
\[
p_{W;t_-,t_+}(\nu)=\sum_{q}\,
\bigl\|E_{W,t_+}(\{q-\nu\})\,E_{W,-t_-}(\{q\})\,\Psi\bigr\|^{2}.
\]
\end{definition}

Note what the protocol does \emph{not} assume: nothing about reflected or
transmitted waves, nothing about how many kinds of excitation exist, nothing
about bound states.  It is a pair of measurements and bookkeeping.

The single dynamical assumption of the main theorem says that these
statistics settle down at late times and do not leak to infinity as the
window grows.  ``Ces\`aro convergence'' below means convergence of running
time-averages: a function $g(t)$ Ces\`aro-converges if
$\frac{1}{T}\int_{T}^{2T}g(t)\,dt$ converges as $T\to\infty$; this is weaker
than convergence of $g(t)$ itself and tolerates persistent oscillations.

\begin{hypothesis}[local relaxation, (LR)]\label{hyp:lr}
For the fixed kink-plus-packet state $\Psi$, the fixed cut $c$, and a fixed
growing family of windows $W_m\uparrow\Z$ each containing $c$:
\begin{enumerate}
\item[(i)] \emph{(settling)} there is one sequence of times $T_n\to\infty$
along which, for every fixed window $W$, both the single-time expectation
values of all observables in $W$ (at times $\pm t$, Ces\`aro-averaged) and
the two-measurement outcome laws $p_{W;t_-,t_+}$ (Ces\`aro-averaged over both
times) converge; call the limiting law $p_W$;
\item[(ii)] \emph{(vanishing back-action on the mean)} the first measurement
does not shift the late-time mean charge in the limit: the Ces\`aro average
of $\bigl\langle \mathcal{D}_{W,-t_-}\bigl(\alpha_{t_+}(\widehat Q_{W,c})\bigr)
-\alpha_{t_+}(\widehat Q_{W,c})\bigr\rangle_\Psi$ tends to zero for every
fixed $W$, where $\mathcal{D}_{W,-t_-}(A):=\sum_q E_{W,-t_-}(\{q\})\,A\,
E_{W,-t_-}(\{q\})$ is the operation of performing the first measurement and
forgetting its outcome.  Only this averaged, first-moment statement is
assumed --- not that the operators commute;
\item[(iii)] \emph{(no charge stranded at the window edge)} the limiting laws
are tight with first moments, uniformly in the window:
\[
\lim_{M\to\infty}\ \sup_{m}\ \sum_{|\nu|>M}\bigl(1+|\nu|\bigr)\,p_{W_m}(\nu)\;=\;0 .
\]
\end{enumerate}
Optionally, one may also assume the laws $p_{W_m}$ converge (weakly, i.e.\
probability by probability) to a law $p$; this optional clause buys only the
uniqueness of the limit values and nothing in the quantization statements.
\end{hypothesis}

\begin{remark}[status of (LR)]
Hypothesis~\ref{hyp:lr} is \emph{assumed}, not derived from the microscopic
dynamics; this is stated prominently in Section~\ref{sec:notproved}.  Its
role is to be the \emph{weakest} sufficient condition found under which the
quantization argument closes.  It is strictly weaker than any scattering
assumption: it does not name channels, does not require wave operators, and
tolerates bound states, absorption, and any number of excitation species.
It can fail --- e.g.\ if charge accumulates permanently at the window
boundary --- and the numerical probe of Section~\ref{sec:numerics} was
designed to look for exactly such failures (none was found).
\end{remark}

\section{The quantization theorems}\label{sec:main}

\begin{theorem}[finite-window lattice rigidity; claim \texttt{M-INDEX-fin}]
\label{thm:fin}
Assume Hypothesis~\ref{hyp:int} and Definition~\ref{def:vacua}(3).  Fix any
finite window $W=[a,b]$ and cut $c\in W$.  Then:
\begin{enumerate}
\item all eigenvalues of $\widehat{Q}_{W,c}$ lie in a single shifted integer
lattice $\kappa_{W,c}+\Z$, where the shift
$\kappa_{W,c}\equiv|W|\kappa+\rho\,(a+b-1-2c) \pmod{\Z}$ depends on the
window, the cut, and the density, but \emph{not} on the state and \emph{not}
on time (time evolution preserves the spectrum of an observable);
\item consequently, in the protocol of Definition~\ref{def:tpm}, both
recorded outcomes $q_\pm$ lie in the \emph{same} shifted lattice, the shifts
cancel in the difference, and every recorded escaped charge $\nu=q_--q_+$ is
an exact integer --- at every window size, at all times, with no hypotheses
on the dynamics beyond its existence.
\end{enumerate}
\end{theorem}

\begin{proof}[Proof sketch]
Each $S^z_x$ has spectrum in $\kappa+\Z$ (Hypothesis~\ref{hyp:int}); the
window charge is a sum of commuting such operators minus the number
$\rho\,(c-a+1)-\rho\,(b-c)$; adding commuting operators adds their spectra.
The integrality of $\nu$ is then arithmetic --- crucially, arithmetic of two
outcomes of the \emph{same} observable at fixed $W$, so the state-independent
shift cancels exactly.  This is where the protocol earns its keep: it is
immune to the noncommutativity trap of Section~\ref{sec:naive} because no
operator subtraction ever occurs, only subtraction of recorded numbers.
\end{proof}

\begin{theorem}[quantization of the memory outcome law; claim
\texttt{M-INDEX-spec}]\label{thm:spec}
Assume the setting of Sections~\ref{sec:setting}--\ref{sec:wall}
(finite-range charge-conserving dynamics, symmetric vacuum pair with density
$\rho>0$, kink state with summable deviations and first moments),
Hypothesis~\ref{hyp:int}, and Hypothesis~\ref{hyp:lr}.  Then, along the time
sequence provided by (LR):
\begin{enumerate}
\item for every fixed window $W_m$ the limiting outcome law $p_{W_m}$ exists
and is supported on the integers;
\item every limit point $p=\{p_\nu\}_{\nu\in\Z}$ of the family
$\{p_{W_m}\}$ as $m\to\infty$ is a probability distribution supported on
$\Z$: no probability escapes to infinity and none lands off the integer
lattice.  (With the optional convergence clause of (LR) the limit $p$ is
unique.);
\item the asymptotic mean wall displacement, defined as the ordered double
limit (time limits first at fixed window, then $W_m\uparrow\Z$) of the
difference of expectations of the wall-position observable
$X_{W,c}$ of Definition~\ref{def:window}, exists along the same subsequence
and equals
\[
\boxed{\ \delta x\;=\;-\,\frac{1}{2\rho}\ \sum_{\nu\in\Z}\ \nu\,p_\nu\ }
\]
\item in particular, \emph{every individual outcome} of the protocol
displaces the recorded wall position by an exact integer multiple of
$1/(2\rho)$ lattice sites, while the \emph{mean} $\delta x$ is a
probability-weighted average of these quantized outcomes and is not itself
quantized.
\end{enumerate}
Bound states, absorption of the magnon by the wall, additional excitation
species, and any other exotic process are all permitted: they contribute
probability at other integers $\nu$ and cannot produce non-integer outcomes.
\end{theorem}

\begin{proof}[Proof sketch]
Clause (1): Theorem~\ref{thm:fin}(2) puts each finite-time law on $\Z$;
Ces\`aro limits of probability laws on the closed set $\Z$ stay on $\Z$
provided no mass escapes, which at fixed $W$ is automatic (finitely many
eigenvalues).  Clause (2): as $W$ grows the lattice itself is fixed
(integers), and Hypothesis~\ref{hyp:lr}(iii) is exactly the tightness needed
so that weak limit points remain probability laws on $\Z$ (this is the
standard compactness argument for tight families, and the counterexamples
showing that plain convergence of means would \emph{not} suffice are recorded
alongside the proof).  Clause (3): at fixed window and finite times, charge
conservation (Definition~\ref{def:dynamics}) gives an exact identity: the
change of $\langle\widehat Q_{W,c}\rangle$ equals the net current through the
two ends, and the first-moment identity
$\sum_\nu\nu\,p_{W;t_-,t_+}(\nu)
=\langle\alpha_{-t_-}(\widehat Q)\rangle-
\langle\mathcal{D}_{W,-t_-}(\alpha_{t_+}(\widehat Q))\rangle$
relates the protocol mean to the unmeasured evolution up to the back-action
term that Hypothesis~\ref{hyp:lr}(ii) kills in the limit;
Hypothesis~\ref{hyp:lr}(iii) (first moments, not just tightness) lets the
mean pass through the window limit.  Dividing by $-2\rho$ converts escaped
charge to wall displacement (Definition~\ref{def:window}).  The full
numbered-step proof is \texttt{theory/memory-index.md}; the first-moment
identity and the tightness counterexamples carry machine certificates.
\end{proof}

\begin{remark}[consistency with the scattering picture]\label{rem:consistency}
Suppose one \emph{additionally} assumes the full scattering package that
earlier work had to assume from the start: a complete two-channel scattering
theory in which the incoming magnon (charge $-1$ relative to the left vacuum)
either reflects (outgoing charge $-1$ on the same side, escaped charge
$\nu=0$) or transmits (outgoing charge $+1$ relative to the right vacuum,
escaped charge $\nu=2$), with transmission probability $T$.  Then the outcome
law of Theorem~\ref{thm:spec} collapses to $p_0=1-T$, $p_2=T$, and the boxed
formula gives $\delta x=-T/\rho$: exactly the previously known conditional
result.  The logical relation is strict: the old theorem = the new theorem +
the scattering package; the quantization itself never needed the package.
\end{remark}

\section{Existence of scattering states from the excitation ansatz}
\label{sec:acex}

Theorem~\ref{thm:spec} quantizes the outcomes without knowing their
probabilities.  Computing the probabilities $p_\nu$ is the job of scattering
theory, and the obstruction there is well known: proving that wave operators
exist and are complete for a many-body lattice system is out of reach in
general, and completeness can even be false as posed (extra channels, bound
states).  The final theorem shows that the \emph{existence} half --- enough
to define the transmission probability on the states the construction
reaches --- follows from hypotheses that the matrix-product
\emph{excitation ansatz} \cite{haegeman2011,haegeman2013,vanderstraeten2019}
is designed to provide.  The philosophy: the ansatz supplies the kinematic
data (which single excitations exist and how they disperse); once those data
are assumed exact, the dynamical statement (scattering states exist) becomes
a theorem.  Completeness is \emph{not} claimed, and by
Theorem~\ref{thm:spec} it is no longer needed for quantization.

\begin{definition}[wave operators]\label{def:wave}
Let $\Hs$ carry the dynamics $e^{-iHt}$ of the kink sector (the
Hilbert-space realization of Definition~\ref{def:gns} for kink states), and
let $\Hs_{\mathrm{ch}}$ be a \emph{channel space}: square-summable functions
of a kink momentum $p$ and a magnon momentum $k$, evolving under
multiplication by $e^{-i(E_K(p)+\omega(k))t}$, where $E_K$ and $\omega$ are
given functions (the \emph{dispersions}: the energy of each excitation as a
function of its momentum; the derivative of a dispersion is the excitation's
propagation speed, its \emph{group velocity}).  Given a linear \emph{identification
map} $J:\Hs_{\mathrm{ch}}\to\Hs$ (built below from the ansatz data), the
\emph{wave operators} are the strong limits
\[
W_\pm\;:=\;\slim_{t\to\pm\infty}\;e^{iHt}\,J\,e^{-iH_{\mathrm{ch}}t}
\]
whenever they exist on a stated domain.  Physically: $W_-\phi$ is the true
state of the interacting chain that, in the far past, looked like the freely
moving kink-plus-magnon configuration $\phi$.
\end{definition}

\begin{hypothesis}[exact band data; abbreviated statement]\label{hyp:band}
There exist: (a) a family of \emph{kink vectors} $\kappa_a(x)\in\Hs$
($x\in\Z$, finitely many labels $a$), translation-covariant, whose closed
span is invariant under $H$ with matrix dispersion $E_K$, twice continuously
differentiable in momentum; (b) over each vacuum, \emph{magnon creation
operators}: almost-local operators (localized up to tails decaying faster
than any power) adding one unit of charge, whose action on the vacuum
realizes an isolated single-magnon energy band with dispersion $\omega$ and
positive gap; both band relations hold \emph{exactly}, not approximately ---
this is the mathematical idealization of ``the excitation ansatz is exact'';
(c) a \emph{two-cluster factorization estimate}: there are $C<\infty$ and
$\lambda\in(0,1)$ such that for all kink vectors and all observables $A$
localized at least $r$ sites to the left and $B$ at least $r$ sites to the
right of the kink positions,
\[
\bigl|\langle\kappa_a(x),\,AB\,\kappa_{a'}(x')\rangle
-\omega_\alpha(A)\,\omega_\beta(B)\,\langle\kappa_a(x),\kappa_{a'}(x')\rangle\bigr|
\;\le\;C\,\|A\|\,\|B\|\,\lambda^{r};
\]
(d) the incoming wave packets are smooth with disjoint group-velocity
supports separated by some $\varepsilon_v>0$ (in particular, this is a
fixed-packet theorem: the zero-velocity limit $k\to0$ is excluded, and no
uniformity in it is claimed).
\end{hypothesis}

\begin{theorem}[existence of kink--magnon scattering states; claim
\texttt{AC-EX}, parts ACE.1--ACE.2]\label{thm:acex}
Under Definition~\ref{def:vacua}, Definition~\ref{def:dynamics} and
Hypothesis~\ref{hyp:band}, with the identification map $J$ built by applying
a momentum-filtered magnon creation operator to a kink vector, the wave
operators $W_-$ (one incoming configuration) and $W_+$ (two outgoing
configurations, magnon on the left or on the right of the kink) exist on the
stated packet domains and are isometries (norm-preserving maps, so no
probability is lost in the construction).  On the \emph{constructed}
outgoing subspace, the orthogonal projection $N_T^{\mathrm{ex}}$ onto the
transmitted configuration is well defined, and
$\langle\Psi,N_T^{\mathrm{ex}}\Psi\rangle$ is the transmission probability of
the constructed channels.  Nothing is asserted about the orthogonal
complement of the constructed subspace: completeness --- the statement that
the constructed states exhaust everything that propagates --- is not claimed,
is not used anywhere in this report, and is known to be unprovable at this
generality.
\end{theorem}

\begin{proof}[Proof sketch]
The standard route (Cook's method): show that the time derivative of
$e^{iHt}J e^{-iH_{\mathrm{ch}}t}\phi$ is integrable in $t$.  The derivative
is controlled by (i) the exactness of the band relations, which makes the
naive error term vanish identically for a single excitation, leaving only
cross terms between the kink and the magnon; (ii) the cluster estimate
(c) and the velocity separation (d), which confine the cross terms to an
exponentially suppressed spatial cone, giving an integrable
$O(|t|^{-3})$ bound.  A technical device is essential and easy to state: the
magnon creation operators must be \emph{momentum-filtered} (their momentum
content cut off smoothly around the packet support) --- with the naive
unfiltered normalization the hopping kernels decay too slowly and the
integral does not close.  The port of this argument from the
translation-invariant setting of \cite{bachmann2015} to the kink sector
(where the wall breaks translation invariance of the \emph{state} but the
sector as a whole still carries a translation action) is the substance of
the proof; every step is re-derived rather than cited by analogy.  Full
proof: \texttt{theory/ansatz-scattering.md}.
\end{proof}

\begin{remark}[bridge to the outcome law; conditional]\label{rem:bridge}
On states in the range of both wave operators, one expects the outcome law of
Theorem~\ref{thm:spec} to be carried by the constructed channels with
$p_2=\langle N_T^{\mathrm{ex}}\rangle$ (as in
Remark~\ref{rem:consistency}).  This bridge is proved \emph{conditionally}
on three explicitly stated extra assumptions: the charge bookkeeping of the
two legs ($q$-values $-1$, $-1$, $+1$ as in Remark~\ref{rem:consistency}), a
local-decay hypothesis (the outgoing state, restricted to any fixed window
around the wall, becomes charge-definite --- the missing lemma is named and
tracked), and the applicability of Hypothesis~\ref{hyp:lr} to the vector in
question.  None of the three is proved here; see
Section~\ref{sec:notproved}.
\end{remark}

\section{Numerical verification}\label{sec:numerics}

Three independent machine checks accompany the theorems.  Two conventions
first.  \emph{Exact diagonalization} (ED) means: represent the full
many-body Hilbert space of a finite chain on a computer with no
approximation beyond machine arithmetic, and evolve by the exact matrix
exponential.  A \emph{red test} (used throughout the project) means: every
machine-checked certificate must be demonstrated to \emph{fail} when one of
its hypotheses is deliberately broken; a check that cannot fail certifies
nothing.  All checkers pass their green runs and fail their red runs; the
suite is rerun before every status change.

\paragraph{Pre-registered falsification attempt.}
Before the proofs of Section~\ref{sec:main} existed, a falsifier was
specified: if the outcome distribution of the window charge, in an honest
simulation of a kink--magnon collision, put persistent probability off the
integer lattice or leaked mass under growing windows, the theorem targets
would have been withdrawn.  The model: the spin-$\tfrac12$ easy-axis XXZ
ferromagnet,
\[
H=-J\sum_x\Bigl[S^x_xS^x_{x+1}+S^y_xS^y_{x+1}
+\Delta\bigl(S^z_xS^z_{x+1}-\tfrac14\bigr)\Bigr],
\qquad \Delta=2,\ J=1,
\]
with boundary spins frozen up (left) and down (right) to hold one wall in
the system; here $\rho=\tfrac12$, so the predicted displacement quantum is
$1/(2\rho)=1$ site per unit of escaped charge, and the transmission outcome
($\nu=2$) displaces by $2$ sites.  A chain of $N=50$ sites (restricted to at
most five domain walls, a truncation validated against the exact space at
$N=22$ to $3.4\times10^{-3}$ in every probability) was prepared with a
\emph{dressed} kink (the true ground state of the kink sector, not a sharp
product wall) and an incoming magnon packet, and evolved through the
collision.  Results: the probability off the integer lattice was
\emph{exactly zero} to machine precision at every time, window, and
geometry tested; at the latest clean time the law concentrated on the two
scattering outcomes, $p(\nu=2)=0.9996$, $p(\nu=0)=0.0004$ (this model is
nearly transparent at the chosen momentum), with residual mass
$\le 2.5\times10^{-3}$ on neighboring integers; the mean displacement
computed from the outcome law agreed with an independent estimator of the
wall position (a profile centroid, algebraically unrelated to the window
charge) to $0.011$ sites.  Two instructive artifacts were found and
documented: preparing a \emph{sharp} wall instead of the dressed one fakes a
$5\%$ off-outcome signal (the preparation defect radiates slow magnon pairs
--- a state-preparation effect, not a violation); and the residual
integer-valued mass sits in near-threshold two-magnon configurations with
group velocity near zero --- precisely the slow-excitation regime that
scattering-based approaches cannot control, absorbed here by
Theorem~\ref{thm:spec} because such outcomes still land on integers.

\paragraph{Certificates for the proofs.}
A second program certifies, in exact or high-precision arithmetic, the
load-bearing computations of Sections~\ref{sec:mps}--\ref{sec:main}: the
lattice arithmetic of Theorem~\ref{thm:fin} including an irrational-offset
stress test; the $2\times2$ counterexample of Section~\ref{sec:naive}; the
divergent-fluctuation counterexample of Theorem~\ref{thm:noq}; the
density-equals-phase-slope computation behind Theorem~\ref{thm:density} on
randomly drawn symmetric tensors of bond dimensions up to $4$ (agreement to
$9\times10^{-16}$, with the necessity of the antisymmetric vacuum pair
red-certified); the first-moment identity of Theorem~\ref{thm:spec}; and the
exact reduction of Remark~\ref{rem:consistency}.  A third program certifies
the two quantitative pillars of Theorem~\ref{thm:acex}: the band and
threshold data of the reference model, and the decay rate of the Cook
integrand, with a red mode confirming that a slowly decaying (unfiltered)
kernel destroys the required integrability --- the failure mode the
momentum filter exists to prevent.

\section{What is not proved}\label{sec:notproved}

For clarity, the complete list of assumptions and open items.

\begin{enumerate}
\item \emph{Hypothesis (LR) is not derived from the dynamics.}  No proof is
offered that any concrete Hamiltonian satisfies the relaxation clauses; (LR)
is the stated interface between the quantization arithmetic and the hard
open problems of many-body relaxation theory.  In particular, no claim of
the form ``scattering assumptions imply (LR)'' is made; an attempted proof
of such an implication was reviewed and rejected during verification.
\item \emph{Hypothesis \ref{hyp:band} (exact band data) is not derived.}
In practice the excitation ansatz provides these data as the output of a
numerical optimization; here their exactness is an assumption.  The cluster estimate (c) is the
load-bearing clause and is verified on no model.
\item \emph{Completeness of scattering is neither proved nor assumed.}  It
is provably not derivable from band data alone, and Theorem~\ref{thm:spec}
was engineered precisely so that it is not needed.
\item \emph{The bridge of Remark~\ref{rem:bridge}} rests on a named,
unproved local-decay lemma and on (LR); both are tracked as open problems.
\item \emph{The mean $\delta x$ is not quantized}, and no claim to the
contrary is made anywhere; quantization is a property of individual
protocol outcomes.
\item \emph{Conjecture (charge operator on the kink representation).}  While
Theorem~\ref{thm:noq} kills the state-independent construction, it is
conjectured that on the specific Hilbert-space realization of a fixed kink
state the symmetry rotations are implemented by a continuous unitary group
whose generator has purely discrete, integer-spaced spectrum; supporting
evidence (uniformly bounded charge variance on such realizations) is
recorded, but no proof exists.
\item \emph{Model scope of the numerics.}  The simulations concern one
ferromagnetic model family; they are consistency checks and falsification
attempts, not proofs, and are labelled as such throughout.
\end{enumerate}

\appendix

\section{How these results were verified}\label{app:method}

Each theorem above was produced by an adversarial loop: a \emph{proposer}
(an automated prover) writes the full argument in numbered, individually
addressable steps, every step justified by a numbered definition, a
previously established claim, or a named machine-checked computation; an
independent \emph{critic} (a different automated system, from a different
model family, to decorrelate blind spots) attacks the argument by
recomputing its key steps, constructing counterexamples, and auditing every
quantifier, and issues a written verdict listing objections classified by
severity, each with a demanded fix and a statement of what survives if the
fix fails.  The loop repeats until a verdict contains no fatal or major
objection.  For the results reported here the loop ran three rounds; the
objection counts fell from $20$ (first round, three parallel verdicts) to
$14$ to $13$, with major objections $6\to3\to1\to0$ and no fatal objection
at any round.  Two of the corrections flowed in the reverse direction ---
errors in a critic's verdict caught and refuted by the machine-checked
certificates --- which is the intended behavior of the design.  Two draft
claims did not survive: the global charge operator (now the refutation,
Theorem~\ref{thm:noq}) and an implication from scattering assumptions to
(LR) (withdrawn; item 1 of Section~\ref{sec:notproved}).  The verdicts,
proofs, and certificates are in the project repository
(\texttt{theory/memory-index.md}, \texttt{theory/ansatz-scattering.md},
\texttt{theory/verdicts/memory-index-r1}--\texttt{r3},
\texttt{theory/checks/}).

\begin{thebibliography}{9}
\bibitem{perez2008} D.~P\'erez-Garc\'ia, M.~M.~Wolf, M.~Sanz, F.~Verstraete,
and J.~I.~Cirac, \emph{String order and symmetries in quantum spin lattices},
Phys.\ Rev.\ Lett.\ \textbf{100}, 167202 (2008); arXiv:0802.0447.
\bibitem{cirac2021} J.~I.~Cirac, D.~P\'erez-Garc\'ia, N.~Schuch, and
F.~Verstraete, \emph{Matrix product states and projected entangled pair
states: concepts, symmetries, theorems}, Rev.\ Mod.\ Phys.\ \textbf{93},
045003 (2021); arXiv:2011.12127.
\bibitem{haegeman2011} J.~Haegeman, B.~Pirvu, D.~J.~Weir, J.~I.~Cirac,
T.~J.~Osborne, H.~Verschelde, and F.~Verstraete, \emph{Variational matrix
product ansatz for dispersion relations}, Phys.\ Rev.\ B \textbf{85}, 100408
(2012); arXiv:1103.2286.
\bibitem{haegeman2013} J.~Haegeman, S.~Michalakis, B.~Nachtergaele,
T.~J.~Osborne, N.~Schuch, and F.~Verstraete, \emph{Elementary excitations in
gapped quantum spin systems}, Phys.\ Rev.\ Lett.\ \textbf{111}, 080401
(2013); arXiv:1305.2176.
\bibitem{vanderstraeten2019} L.~Vanderstraeten, J.~Haegeman, and
F.~Verstraete, \emph{Tangent-space methods for uniform matrix product
states}, SciPost Phys.\ Lect.\ Notes \textbf{7} (2019); arXiv:1810.07006.
\bibitem{bachmann2015} S.~Bachmann, W.~Dybalski, and P.~Naaijkens,
\emph{Lieb--Robinson bounds, Arveson spectrum and Haag--Ruelle scattering
theory for gapped quantum spin systems}, Ann.\ H.\ Poincar\'e \textbf{17},
1737 (2016); arXiv:1412.2970.
\bibitem{bbdf2020} S.~Bachmann, A.~Bols, W.~De Roeck, and M.~Fraas,
\emph{A many-body index for quantum charge transport}, Commun.\ Math.\
Phys.\ \textbf{375}, 1249 (2020); arXiv:1810.07351.
\end{thebibliography}

\end{document}
